%% TENSOR-BASED COGNITIVE OVERSIGHT (TCO)
%% ACM submission draft — April 2026
%% Compile: pdflatex main.tex → bibtex main → pdflatex main.tex (×2)
%%
%% Target venues:
%%   CHI 2027   → \documentclass[sigconf,anonymous]{acmart}
%%   EMSE       → \documentclass[acmlarge]{acmart}
%%   IEEE RE    → switch to IEEEtran.cls
%%
\documentclass[manuscript,review,anonymous]{acmart}

\usepackage{booktabs}
\usepackage{amsmath}
\usepackage{amssymb}
\usepackage{listings}
\usepackage{xcolor}
\usepackage{multirow}
\usepackage{tabularx}

%% --- listings style for Python pseudocode ---
\lstset{
  language=Python,
  basicstyle=\ttfamily\small,
  keywordstyle=\color{blue}\bfseries,
  commentstyle=\color{gray}\itshape,
  stringstyle=\color{green!60!black},
  showstringspaces=false,
  breaklines=true,
  frame=single,
  rulecolor=\color{black!20},
  numbers=left,
  numberstyle=\tiny\color{gray},
  xleftmargin=1.5em,
}

%% --- ACM metadata ---
\acmConference[CHI '27]{ACM CHI Conference on Human Factors in Computing Systems}
  {April 26--May 1, 2027}{Yokohama, Japan}
\acmBooktitle{Proceedings of the 2027 CHI Conference on Human Factors in Computing Systems}
\acmISBN{978-1-4503-XXXX-X/27/04}
\acmDOI{10.1145/XXXXXXX.XXXXXXX}
\copyrightyear{2027}
\acmYear{2027}

\begin{document}

\title{Tensor-Based Cognitive Oversight (TCO): A Framework for Human Orchestration of Multi-Agent AI Software Systems}

\author{Juan Pablo Chancay}
\affiliation{%
  \institution{Aural Syncro}
  \city{[City]}
  \country{[Country]}
}
\email{[email]}

%%
%% Abstract
%%
\begin{abstract}
Multi-agent AI systems generating software artifacts at continuous scale create a
fundamental bottleneck: human supervisors must review outputs that systematically
exceed their cognitive processing capacity. Existing Human-in-the-Loop (HITL)
frameworks address this at the artifact level, leaving review effort coupled linearly
to output volume. We propose \textit{Tensor-Based Cognitive Oversight (TCO)}, a
framework that restructures human oversight by (1)~transforming artifacts into
normalized 11-dimensional evaluation vectors \(\mathbf{V} \in [0,1]^{11}\), (2)~aggregating
vectors into a cognitive tensor \(T[d, i, j, k] \in \mathbb{R}^{n \times s \times a \times t}\)
capturing quality state across dimensions, stages, agents, and time, (3)~deriving
actionable inferences \(I: T \rightarrow \{\Omega, \Delta, \mathcal{P}, \Xi\}\), and
(4)~enabling bidirectional policy injection in natural language. The framework introduces
the concept of the \textit{Natural Cognitive Frontier (NCF)}: the abstraction level at
which cognitive demand is calibrated to human working memory capacity. We present an
experimental design (between-subjects, $n = 40$, 5 fault-injection scenarios) to test
five hypotheses: reduced cognitive load, improved decision accuracy, earlier fault
detection, lower time-to-correct, and higher policy injection quality. We report
the TCO architecture, a deployable MVP implementation, and implications for human
oversight of autonomous AI development pipelines.
\end{abstract}

\keywords{human-in-the-loop, cognitive oversight, multi-agent AI, software quality,
  evaluation vector, natural cognitive frontier, human-AI collaboration}

\begin{CCSXML}
<ccs2012>
   <concept>
       <concept_id>10003120.10003121.10003122</concept_id>
       <concept_desc>Human-centered computing~HCI theory, concepts and models</concept_desc>
       <concept_significance>500</concept_significance>
   </concept>
   <concept>
       <concept_id>10011007.10011074.10011092</concept_id>
       <concept_desc>Software and its engineering~Software quality</concept_desc>
       <concept_significance>300</concept_significance>
   </concept>
   <concept>
       <concept_id>10010147.10010178.10010179</concept_id>
       <concept_desc>Computing methodologies~Intelligent agents</concept_desc>
       <concept_significance>300</concept_significance>
   </concept>
</ccs2012>
\end{CCSXML}
\ccsdesc[500]{Human-centered computing~HCI theory, concepts and models}
\ccsdesc[300]{Software and its engineering~Software quality}
\ccsdesc[300]{Computing methodologies~Intelligent agents}

\maketitle

%%
%% 1. Introduction
%%
\section{Introduction}

Artificial intelligence systems are increasingly capable of generating complete software
artifacts at scale --- code, configurations, architectural decisions, and deployment
pipelines. While this capability enhances productivity, it introduces a structural
bottleneck: human supervisors are required to validate outputs that systematically
exceed efficient cognitive processing capacity \cite{sweller1988}.

This bottleneck manifests as cognitive saturation. When working memory is overloaded,
error detection deteriorates, decision quality degrades, and human agents shift from
deliberate analytical mode to reactive heuristic mode \cite{sweller2011}. Cowan
\cite{cowan2001} establishes a working memory capacity of 3--5 meaningful chunks;
multi-agent software pipelines routinely produce dozens of interdependent artifacts per
cycle, each requiring cross-referencing judgment.

Traditional HITL models assume bounded, interpretable discrete outputs at well-defined
checkpoints \cite{hitlreview2025}. In continuous multi-agent AI-driven development,
this assumption fails: outputs are continuous, each artifact is complex, and errors
propagate across interdependent components. Recent industry analysis confirms the
consequence: Codebridge's study of 211 million changed lines of code over four years
\cite{codebridge2025} documents increasing defect rates in repositories with high
AI-generated code proportions.

This paper argues that the bottleneck is not human capability but the \textbf{level of
abstraction} at which humans interact with AI outputs. We propose
\textit{Tensor-Based Cognitive Oversight (TCO)}, a framework that replaces
per-artifact review with system-level orchestration at the
\textit{Natural Cognitive Frontier (NCF)} --- the abstraction level where cognitive
demand is calibrated to human working memory.

TCO makes four technical contributions:
\begin{enumerate}
  \item A formal \textbf{evaluation vector model} $\mathbf{V} \in [0,1]^{11}$ normalising
    artifact quality across 11 semantic dimensions.
  \item A \textbf{cognitive tensor} $T[d,i,j,k]$ representing global system state
    across dimensions, pipeline stages, agents, and time.
  \item An \textbf{inference model} $I: T \rightarrow \{\Omega, \Delta, \mathcal{P}, \Xi\}$
    deriving actionable system state, trends, conflicts, and recommendations.
  \item A \textbf{bidirectional policy loop} enabling natural language re-orchestration
    of the agent system from the NCF.
\end{enumerate}

%%
%% 2. Related Work
%%
\section{Related Work}

\subsection{Human-in-the-Loop Frameworks for Agentic AI}

A 2025 systematic review of over 400 HITL implementations identifies cognitive load
and trust calibration as persistently unsolved challenges in deployed systems
\cite{hitlreview2025}. HITL effectiveness degrades sharply as output volume and
artifact complexity increase.

The most directly related work is HULA \cite{hula2024}, which integrates human review
checkpoints into a multi-agent software development pipeline. TCO and HULA share the
recognition that unmediated agentic output is unmanageable; they differ architecturally.
HULA preserves artifact-level intervention --- review effort scales with volume.
TCO proposes system-level orchestration, decoupling cognitive effort from output volume.
Pareschi and Saghafian \cite{pareschi2024} argue for a shift from direct to delegated
control as agent capability increases; TCO operationalizes that delegation with the
vector $\phi$, tensor $T$, and policy loop.

\subsection{Supervisory Control and Levels of Automation}

Parasuraman, Sheridan, and Wickens \cite{parasuraman2000levels} establish a ten-level
taxonomy of human-automation interaction, predicting that as AI autonomy increases,
supervisory interfaces must shift from monitoring individual actions to managing
goal-level outcomes. Miller and Parasuraman \cite{miller2007} show that effective
oversight requires interfaces calibrated to the decision horizon, not operational detail.
Cummings \cite{cummings2014} demonstrates empirically that performance bottlenecks in
complex automation arise from interface design failures: supervisors consistently
perform better with aggregated decision-relevant displays than with technical readouts.

\subsection{Observability and Monitoring Tools}

Commercial platforms (Datadog, New Relic, Grafana) excel at infrastructure-level
telemetry but cannot represent semantic quality dimensions
(architectural alignment, security exposure, maintainability trajectory) that require
natural language understanding. The NIST AI Risk Management Framework \cite{nist2023rmf}
mandates continuous monitoring but does not address the cognitive interface problem.
Sculley et al.\ \cite{sculley2015} document that technical debt in ML systems
accumulates at the semantic level before manifesting in operational metrics; TCO's
tensor is designed to surface this pre-operational degradation.

\subsection{Cognitive Load in Software Engineering}

Code review is among the most cognitively demanding activities in software development,
with quality degradation documented beyond 60--90 minutes of sustained review
\cite{sweller2011}. Amershi et al.\ \cite{amershi2019} document compound cognitive
overhead when ML components are involved. Codebridge \cite{codebridge2025} provides
large-scale evidence that the binding constraint in AI-assisted development is review
throughput, not generation throughput.

\subsection{Multi-Agent Evaluation and LLM-as-Judge}

LLM-based systems --- SWE-agent \cite{sweagent2024}, Devin, OpenHands --- generate
semantically interdependent artifacts across stages. TCO's tensor indexing
$T[d, i, j, k]$ explicitly captures cross-agent dependency: the conflict component
$\mathcal{P}$ detects inter-agent disagreement before downstream failures occur.
The LLM-as-Judge paradigm \cite{llmjudge2025winlp,llminterrater2025,llmjudgeempirical2025}
provides the substrate for automated vector generation, with inter-rater reliability
$\kappa \geq 0.70$ when structured rubrics are used \cite{llminterrater2025}.

%%
%% 3. Theoretical Framework
%%
\section{Theoretical Framework}

TCO is grounded in four established theories, integrated within the specific domain of
human orchestration of multi-agent AI software systems.

\subsection{Cognitive Load Theory}

Sweller's Cognitive Load Theory \cite{sweller1988} distinguishes intrinsic,
extraneous, and germane load competing for finite working memory capacity
(3--5 chunks \cite{cowan2001}). Artifact review imposes maximum intrinsic load
(complex outputs) and maximum extraneous load (machine-optimised format). TCO
eliminates extraneous load via the vector $\phi(A)$ and intrinsic load via the tensor
$T$, leaving the human with pure germane load: interpretation, judgment, and policy.

\subsection{Situation Awareness}

Endsley's SA model \cite{endsley1995} specifies three levels --- L1 (perception),
L2 (comprehension), L3 (projection) --- with 76\% of SA errors occurring at
L1 \cite{endsley2000}. HITL systems force operation at L1 (reading code, parsing logs).
TCO shifts human operation to L2/L3: the vectorisation layer handles L1 automatically.

\begin{table}[t]
  \caption{Situation Awareness Levels Mapped to TCO Architecture}
  \label{tab:sa-mapping}
  \begin{tabular}{p{3cm}p{5cm}}
    \toprule
    SA Level & TCO Component \\
    \midrule
    L1 --- Perception & Vectorisation layer $\phi$: automated extraction \\
    L2 --- Comprehension & Tensor $T$ + Inference layer: aggregated state \\
    L3 --- Projection & Trend analysis $\Delta$ and conflict detection $\mathcal{P}$ \\
    Decision \& Action & NCF: natural language policy formulation \\
    \bottomrule
  \end{tabular}
\end{table}

\subsection{Supervisory Control Theory}

Sheridan \cite{sheridan1992} describes the evolution from direct operators to
supervisors in complex automated systems. The aviation shift to glass cockpits
provides the historical precedent: aggregated system-state displays replaced
direct instrument reading, enabling pilots to operate at goal-oriented levels.
TCO proposes the equivalent shift for AI-driven software.

\subsection{Gap in Existing Literature}

Existing approaches fail on two dimensions: (1) they treat HITL as output-level
intervention rather than system-level orchestration; (2) they do not address the
representation problem. TCO is structurally distinct from existing observability
platforms along four dimensions: semantic quality vocabulary, multi-agent aggregation,
cognitive interface orientation, and bidirectional policy loop.

%%
%% 4. Problem Statement
%%
\section{Problem Statement}

\subsection{Cognitive Mismatch in AI Supervision}

Human supervisors of multi-agent AI software systems must simultaneously:
monitor \textit{volume} (dozens of artifacts per cycle), assess
\textit{cross-artifact consistency} (inter-agent coherence), and maintain
\textit{temporal awareness} (quality trends across cycles). Each demand is individually
manageable; their simultaneous imposition saturates working memory.

The root cause is a representation mismatch: AI systems generate outputs in
machine-optimised formats (Python ASTs, YAML trees, deployment manifests), while human
supervisors require decision-oriented abstractions.

\subsection{Incorrect Abstraction Level}

Current HITL systems position the human at the artifact level (Direct control in
Parasuraman's taxonomy \cite{parasuraman2000levels}). The appropriate level for a
human orchestrator of an autonomous pipeline is Goal-Oriented: the human specifies
outcomes, the system reports on trajectory. TCO implements this shift.

\subsection{Automation Bias and TCO's Mitigation}

Automation bias \cite{parasuraman2010,kucking2024} --- over-reliance on automated
outputs --- is a risk at any abstraction level. TCO mitigates this through:
(a) explicit anomaly scoring ($v_{11}$) that surfaces surprising deviations;
(b) inter-agent conflict detection ($\mathcal{P}$) that flags inconsistencies the
human may miss;
(c) confidence scoring ($v_{10}$) that signals low-reliability evaluations requiring
direct review.

%%
%% 5. The TCO Framework
%%
\section{The TCO Framework}

\subsection{Architectural Overview}

TCO implements a six-layer bidirectional architecture:

\begin{enumerate}
  \item \textbf{AI Generation:} Multi-agent pipeline produces artifacts per cycle.
  \item \textbf{QA Evaluation:} LLM-based quality agent evaluates each artifact.
  \item \textbf{Vectorisation $\phi$:} Artifacts $\rightarrow$ $\mathbf{V} \in [0,1]^{11}$.
  \item \textbf{Tensor Aggregation $f$:} $\{\mathbf{V}\} \rightarrow T[d,i,j,k]$.
  \item \textbf{Inference $I$:} $T \rightarrow \{\Omega, \Delta, \mathcal{P}, \Xi\}$.
  \item \textbf{Human Orchestration at NCF:} Natural language policy $\rightarrow$
    re-orchestration of agents.
\end{enumerate}

The bidirectional loop: artifacts flow upward through layers 1--5 to the human;
policy flows downward from the human through layer 6 back to layer 1.

\subsection{The Evaluation Vector Model}

\begin{equation}
  \mathbf{V} = (v_1, v_2, \ldots, v_{11}) \in [0,1]^{11}
\end{equation}

\begin{table}[t]
  \caption{The 11-Dimension Evaluation Vector.
    \textbf{SE} = supervisory estimator (LLM-QA heuristic semantic signal;
    no universal ground truth).
    \textbf{GT} = verifiable ground truth (deterministic static analysis).}
  \label{tab:vector}
  \begin{tabular}{clccc}
    \toprule
    \# & Dimension & Source & Type & Inv. \\
    \midrule
    $v_1$ & functional\_correctness & LLM-QA & \textbf{SE} & --- \\
    $v_2$ & architectural\_alignment & LLM-QA & \textbf{SE} & --- \\
    $v_3$ & scalability\_projection & LLM-QA & \textbf{SE} & --- \\
    $v_4$ & security\_risk & Bandit & GT & $\checkmark$ \\
    $v_5$ & observability\_coverage & Radon & GT & --- \\
    $v_6$ & testability & Radon & GT & $\checkmark$ \\
    $v_7$ & maintainability & Radon & GT & $\checkmark$ \\
    $v_8$ & technical\_debt & Radon & GT & $\checkmark$ \\
    $v_9$ & performance & LLM-QA & \textbf{SE} & --- \\
    $v_{10}$ & confidence & Consensus & --- & --- \\
    $v_{11}$ & anomaly\_score & Z-score & --- & $\checkmark$ \\
    \bottomrule
  \end{tabular}
\end{table}

\noindent \textbf{Inverted dimensions:} $v_4$, $v_8$, $v_{11}$ represent adverse
properties; higher values indicate \textit{worse} quality. The vectorizer normalises
these so that $1.0$ always corresponds to the best quality state.

\medskip
\noindent \textbf{Approximately orthogonal supervisory dimensions.}
The 11 dimensions are \textit{supervisoriamente distinguishable}: each captures a
distinct oversight concern that provides non-redundant decision-relevant signal to the
orchestrator, even when dimensions partially correlate empirically
(e.g.\ $v_7$ maintainability $\leftrightarrow$ $v_8$ technical debt).
An artifact may simultaneously have high $v_1$ (functional correctness) and low $v_7$
(maintainability) --- this is diagnostic information, not contradiction.
Strict statistical independence is not claimed; supervisory distinguishability is
\citep[cf.][for the analogous argument in MQM translation evaluation]{lommel2014}.
Pre-pilot calibration (\S\ref{sec:calibration}) reports pairwise Spearman~$\rho$ between
static-analysis dimensions; any pair with $\rho > 0.90$ receives explicit justification.

\subsection{The Cognitive Tensor}

\begin{equation}
  T[d, i, j, k] \in \mathbb{R}^{n \times s \times a \times t}
\end{equation}

where $n = 11$ (dimensions), $s = 4$ (stages: design, build, test, deploy),
$a$ = number of agents (dynamic), $t$ = cycle index (dynamic).

The tensor enables operations unavailable to per-artifact evaluation: current
snapshot $T[:,:,:,k]$, agent trajectory $T[:,:,j,:]$, stage profile $T[:,i,:,-1]$,
and dimension trend $T[d,:,:,:]$.

\subsection{The Inference Model}

\begin{equation}
  I: T \rightarrow \{\Omega, \Delta, \mathcal{P}, \Xi\}
\end{equation}

\begin{itemize}
  \item $\Omega$ (\textbf{global state}): $\{$stable $(\bar{T} \geq 0.70)$, warning
    $(0.50 \leq \bar{T} < 0.70)$, critical $(\bar{T} < 0.50)\}$
  \item $\Delta$ (\textbf{trend}): $\Delta[d,i,j] = T[d,i,j,k] - T[d,i,j,k-1]$
    for $|\Delta| > 0.05$
  \item $\mathcal{P}$ (\textbf{conflict}): $|T[d,i,j_1,k] - T[d,i,j_2,k]| > 0.30$
  \item $\Xi$ (\textbf{recommendations}): prioritised by estimated impact, derived
    from $\Delta$ and $\mathcal{P}$
\end{itemize}

\subsection{The Natural Cognitive Frontier}

The Natural Cognitive Frontier (NCF) is the level of abstraction at which:
(1) cognitive demand does not exceed working memory capacity;
(2) the human retains sufficient agency to detect novel failure modes;
(3) policy expression is natural (natural language, not code or configuration).

TCO positions the human precisely at the NCF by replacing artifact-level review with
tensor-derived inference output.

\subsection{The Bidirectional Policy Loop}

The human formulates policy as natural language: \textit{``Prioritise security
hardening in the code agent for the next 3 cycles''}. The
\texttt{PolicyProcessor} extracts a structured \texttt{PolicyIntent}:

\begin{lstlisting}[caption={PolicyIntent structured extraction}]
@dataclass
class PolicyIntent:
    target_agents:       list[str]  # ["code_agent"]
    action_type:         str        # "prioritize"
    affected_dimensions: list[str]  # ["security_risk"]
    constraint:          str        # "use parameterized queries"
    priority:            str        # "high"
\end{lstlisting}

\noindent The struct is serialised as a system prompt patch per target agent.
If extraction confidence $< 0.70$, the raw policy text is injected directly
(logged as a degraded case for H5 analysis).

%%
%% 6. Research Objectives and Hypotheses
%%
\section{Research Objectives and Hypotheses}

\noindent \textbf{General Objective:} Determine whether TCO reduces cognitive load and
improves oversight effectiveness in human supervision of multi-agent AI software systems.

\smallskip
\noindent We test five directional hypotheses (between-subjects, $n = 40$,
experimental group: TCO dashboard; control group: raw artifact review):

\begin{description}
  \item[H1] TCO participants report lower NASA Raw-TLX scores than control
    ($p < 0.05$, Mann-Whitney U).
  \item[H2] TCO participants achieve higher fault detection accuracy
    (precision + recall on ground truth faults), with Cohen's $d > 0.50$.
  \item[H3] TCO participants achieve the same detection accuracy on
    temporally accumulated faults (S3 scenario) as on instantaneous faults (S1, S2),
    while control accuracy decreases.
  \item[H4] TCO participants achieve lower time-to-correct on S1--S5 faults
    ($p < 0.05$, Wilcoxon signed-rank).
  \item[H5] TCO participants formulate higher Policy Injection Quality (PIQ)
    structured intents, scored by LLM-Judge against expert rubric
    ($\kappa \geq 0.70$, Spearman $\rho > 0$, $p < 0.05$).
\end{description}

%%
%% 7. Experimental Design
%%
\section{Experimental Design}

\subsection{Overview}

\begin{table}[t]
  \caption{Experimental Design Summary}
  \label{tab:design}
  \begin{tabular}{ll}
    \toprule
    Parameter & Value \\
    \midrule
    Design & Between-subjects, single-blind \\
    $n$ & 40 (20 experimental / 20 control) \\
    Target profile & Software developers, 2+ years experience \\
    Sessions & 1 session per participant, 90--120 min \\
    Tasks & T1--T4 (fault identification, vector analysis, \\
           & early detection, re-orchestration) \\
    Scenarios & S1--S5 (SQL injection, circular dep., tech debt, \\
              & missing metrics, inter-agent conflict) \\
    Measures & NASA Raw-TLX, task accuracy, time-to-correct, PIQ \\
    Analysis & Mann-Whitney U, Cohen's $d$, Spearman $\rho$ \\
    \bottomrule
  \end{tabular}
\end{table}

\subsection{The Five Fault Scenarios}

\begin{table}[t]
  \caption{Fault Scenarios and Ground Truth Vector Deltas}
  \label{tab:scenarios}
  \begin{tabular}{p{0.5cm}p{2.5cm}p{4.5cm}}
    \toprule
    ID & Fault Type & Ground Truth Impact \\
    \midrule
    S1 & SQL injection & $\Delta v_4 = -0.61$, $\Delta v_{11} = -0.59$ \\
    S2 & Circular dependency & $\Delta v_2 = -0.41$, $\Delta v_7 = -0.26$ \\
    S3 & Tech debt accumulation & $\Delta v_8 = -0.08/\text{cycle}$ over 3 cycles \\
    S4 & Missing metrics export & $\Delta v_5 = -0.56$, $\Delta v_9 = -0.26$ \\
    S5 & Inter-agent conflict & $\Delta \mathcal{P}_\rho = +0.41$ \\
    \bottomrule
  \end{tabular}
\end{table}

\subsection{Control Group Environment}

The control group uses the \textit{ControlGroupViewer}: a standardised read-only
multi-tab interface presenting raw Python source, YAML configurations, architecture
markdown, and CI/CD logs with syntax highlighting. A structured correction form
captures free-text issue descriptions with severity classification (Low/Medium/High).
Layout dimensions, font size, and timer display are identical to the TCO dashboard,
eliminating interface ergonomics as a confound.

\subsection{Threats to Validity}

\noindent \textbf{Internal.} Task time-boxing (T1:~25 min, T2:~20 min, T3:~15 min,
T4:~30~min) controls fatigue and equalises exposure. Single-blind design prevents
demand characteristics but risks experimenter bias; mitigated by automated
data collection.

\noindent \textbf{Construct.} QA evaluation circularity risk: the QA agent is both the
measurement instrument (generating vectors) and a component affecting outcomes.
Mitigated by Spearman $\rho \geq 0.75$ validation between LLM scores and static
analysis ground truth on $v_4$, $v_6$, $v_7$, $v_8$ before the pilot.
Cognitive fatigue asymmetry: TCO group may have lower fatigue due to lower cognitive
demand, confounding T4 re-orchestration scores. Mitigated by counterbalanced scenario
ordering and between-subjects design.

\noindent \textbf{External.} Convenience sample; $n = 40$ limits power for small
effects. Controlled deployment scenario; real multi-agent pipelines may exhibit
different fault distributions.

%%
%% 8. Software Architecture
%%
\section{Software Architecture}

\subsection{System Overview}

The TCO MVP consists of four services: (1) a simulated multi-agent pipeline
(LangGraph + Claude API), (2) the TCO Engine (FastAPI + PostgreSQL + Redis),
(3) the orchestration dashboard (React + TypeScript + Recharts), and (4) an
interaction logger capturing millisecond-precision events for all tasks.

\subsection{Evaluation Vector Pipeline}

The vectorizer $\phi$ combines three sources:
\begin{itemize}
  \item \textbf{Radon} (static): cyclomatic complexity ($v_6$), Halstead volume
    ($v_7$), maintainability index ($v_8$).
  \item \textbf{Bandit} (static): CVSS-equivalent severity scoring ($v_4$).
  \item \textbf{LLM-QA} (\textbf{SE} --- supervisory estimators): structured tool-use
    evaluation of $v_1, v_2, v_3, v_9$ using few-shot prompting. These dimensions are
    heuristic semantic signals providing decision-relevant information to the orchestrator;
    they do not constitute objective quality certificates and have no universal ground truth.
\end{itemize}

Confidence ($v_{10}$) is computed as inter-rater agreement between radon and LLM-QA
on the three overlap dimensions (functional correctness, maintainability, testability):

\begin{equation}
  v_{10} = 1 - \frac{1}{3} \sum_{d \in \mathcal{D}_\text{overlap}}
  \left| v_d^{\text{static}} - v_d^{\text{LLM}} \right|
\end{equation}

Anomaly score ($v_{11}$) is the Z-score of current metrics vs historical baseline,
clipped at $3\sigma$ and normalised to $[0,1]$:

\begin{equation}
  v_{11} = 1 - \min\!\left(1, \frac{1}{3} \cdot
  \mathbb{E}\!\left[\left|\frac{x - \mu}{\sigma + \varepsilon}\right|\right]\right)
\end{equation}

\subsection{Artifact Cache}

Identical artifacts (keyed by SHA-256 of content) are not re-evaluated.
Estimated API cost reduction: 40--60\% in repeated-cycle scenarios (DT-017).
Redis TTL: none for artifact hash (deterministic); 30~s for tensor snapshots.

%%
%% 9. Challenges and Open Problems
%%
\section{Challenges and Open Problems}

\paragraph{Vector Calibration.} The 11-dimensional mapping is theoretically grounded
but requires empirical calibration to specific software domains. Spearman $\rho \geq 0.75$
between LLM evaluations and static analysis ground truth is the validation gate.

\paragraph{Tensor Interpretability.} While each cell $T[d,i,j,k]$ is individually
interpretable, the combined tensor is a high-dimensional object. Dimensionality reduction
(PCA, UMAP) for visualisation without losing diagnostic information is an open problem.

\paragraph{Multi-Agent Evaluation Bias.} The QA agent evaluates artifacts generated by
other agents in the same pipeline, creating potential positive correlation bias in scores.
Cross-agent evaluation protocols are an open research direction.

\paragraph{Policy Injection Semantics.} Natural language policies may be ambiguous or
under-constrained. The structured extraction + fallback architecture handles the common
case but does not guarantee semantic fidelity for complex multi-target policies.

%%
%% 10. Expected Contributions
%%
\section{Expected Contributions}

\subsection{Theoretical}

\begin{itemize}
  \item Formal definition of the Natural Cognitive Frontier (NCF) as a design
    construct for human-AI collaboration interfaces.
  \item The cognitive tensor as a representation structure for multi-agent quality state.
  \item The bidirectional policy loop as an operational model for human orchestration
    of autonomous AI pipelines.
\end{itemize}

\subsection{Practical}

\begin{itemize}
  \item Open-source TCO reference implementation (AGPL-3.0): vectorizer,
    aggregator, inference engine, orchestration dashboard.
  \item Deployable experimental platform for replication with other agent
    architectures and software domains.
  \item Empirically validated thresholds for the 11-dimension vector in the
    software development domain.
\end{itemize}

\subsection{Future Work}

\begin{itemize}
  \item Extension to non-software agentic domains (data pipelines, scientific
    workflows, legal document generation).
  \item Adaptive NCF: dynamically adjusting the abstraction level based on
    real-time cognitive load estimates from interaction patterns.
  \item Longitudinal study of operator trust calibration and skill development
    under TCO supervision.
\end{itemize}

%%
%% 11. References
%%
\bibliographystyle{ACM-Reference-Format}
\bibliography{references}

%%
%% Appendix A — Evaluation Vector Schema
%%
\appendix
\section{Evaluation Vector JSON Schema}

\begin{lstlisting}[language={}]
{
  "artifact_id":            "string (UUID)",
  "agent_id":               "string",
  "stage":                  "design|build|test|deploy",
  "cycle_k":                "integer",
  "vector": {
    "functional_correctness":  "float [0,1]",
    "architectural_alignment": "float [0,1]",
    "scalability_projection":  "float [0,1]",
    "security_risk":           "float [0,1]  // inverted",
    "observability_coverage":  "float [0,1]",
    "testability":             "float [0,1]",
    "maintainability":         "float [0,1]",
    "technical_debt":          "float [0,1]  // inverted",
    "performance":             "float [0,1]",
    "confidence":              "float [0,1]",
    "anomaly_score":           "float [0,1]  // inverted"
  },
  "created_at": "ISO-8601 timestamp"
}
\end{lstlisting}

%%
%% Appendix B — Technology Stack
%%
\section{Technology Stack}

\begin{table}[h]
  \caption{MVP Technology Stack}
  \label{tab:stack}
  \begin{tabular}{llp{4cm}}
    \toprule
    Component & Technology & Rationale \\
    \midrule
    LLM & Claude (claude-sonnet-4-6) & Tool-use structured output \\
    Orchestration & LangGraph 0.2+ & State-graph agent coordination \\
    Engine API & FastAPI + Pydantic v2 & High-performance async REST \\
    Storage & PostgreSQL 16 & TimescaleDB upgrade path \\
    Cache & Redis 7 & Artifact + tensor TTL caching \\
    Static analysis & radon + bandit + pylint & No external server required \\
    Dashboard & React 18 + Recharts & Real-time tensor visualisation \\
    \bottomrule
  \end{tabular}
\end{table}

\end{document}
